\documentclass{article}
\usepackage{amsmath}
\usepackage{amssymb}
\usepackage{listings}
\usepackage{xcolor}

% Custom styling for Python code blocks
\definecolor{codegreen}{rgb}{0,0.6,0}
\definecolor{codegray}{rgb}{0.5,0.5,0.5}
\definecolor{codepurple}{rgb}{0.58,0,0.8}
\definecolor{backcolour}{rgb}{0.95,0.95,0.92}

\lstdefinestyle{pythonstyle}{
    backgroundcolor=\color{backcolour},   
    commentstyle=\color{codegreen},
    keywordstyle=\color{magenta},
    numberstyle=\tiny\color{codegray},
    stringstyle=\color{codepurple},
    basicstyle=\ttfamily\footnotesize,
    breakatwhitespace=false,         
    breaklines=true,                 
    captionpos=b,                    
    keepspaces=true,                 
    numbers=left,                    
    numbersep=5pt,                  
    showspaces=false,                
    showstringspaces=false,
    showtabs=false,                  
    tabsize=4
}

\begin{document}

\section*{The Structural Graph Broken Link Problem}

In a custom scalar autograd engine, a Multi-Layer Perceptron (MLP) computes the forward pass by sequentially mapping layers. Let $L^{(k)}$ represent the $k$-th layer of the network.

\subsection*{The Mathematical Graph Breakdown}

For an intermediate hidden layer with $m$ neurons ($m > 1$), the forward transformation yields a vector (or an ordered sequence) of computational graph nodes:
\[
L^{(k)}(x) = \mathbf{a}^{(k)} = \begin{bmatrix} a_1^{(k)} & a_2^{(k)} & \dots & a_m^{(k)} \end{bmatrix}^T \in \mathbb{R}^m
\]
When this vector is passed to the subsequent layer $L^{(k+1)}$, each neuron $j$ in that layer performs a dot product over the iterable collection of inputs:
\[
z_j^{(k+1)} = \sum_{i=1}^{m} w_{ji}^{(k+1)} a_i^{(k)} + b_j^{(k+1)}
\]
During backpropagation, the autograd engine relies on these explicit multiplicative links to propagate the upstream gradient $\frac{\partial \mathcal{L}}{\partial z_j^{(k+1)}}$ backwards via the Chain Rule:
\[
\frac{\partial \mathcal{L}}{\partial a_i^{(k)}} = \sum_{j} \frac{\partial \mathcal{L}}{\partial z_j^{(k+1)}} \cdot w_{ji}^{(k+1)}
\]

\subsection*{The Functional Rupture}

When the data flows into a single-neuron layer (such as a binary classification output layer where $m = 1$), the expected vector collapses to a scalar value:
\[
\mathbf{a}^{(k)} = \begin{bmatrix} a_1^{(k)} \end{bmatrix} \in \mathbb{R}^1
\]
The structural flaw occurs when the code strips the sequence wrapper, changing the data type from an iterable collection to a naked, isolated scalar object:
\[
\mathbf{a}^{(k)} \longrightarrow a_1^{(k)}
\]
If a subsequent processing step or deep architectural layer attempts to consume this output, the initialization mechanism executes a zip operation over the weights and inputs:
\[
\text{zip}\left(\mathbf{w}_j^{(k+1)}, a_1^{(k)}\right)
\]
Because a single scalar \texttt{Value} object lacks an entry dimension ($\text{dim} \notin \mathbb{R}^n$), the iterable matching sequence immediately truncates to an empty set:
\[
\emptyset = \{\}
\]
Consequently, the linear combination evaluates to a dead end:
\[
z_j^{(k+1)} = \sum(\emptyset) + b_j^{(k+1)} = b_j^{(k+1)}
\]
The true input elements $a_i^{(k)}$ are completely omitted from the equation. Because the mathematical multiplication $w \cdot x$ never occurs in memory, no computational dependency edges are created in the autograd directed acyclic graph (DAG).

When calling the backward pass, the partial derivative of the output with respect to the input evaluates to a hard zero:
\[
\frac{\partial z_j^{(k+1)}}{\partial a_i^{(k)}} = 0
\]
The gradient chain vanishes completely at this junction. The backpropagation signal cannot bridge the missing link, leaving all previous layer weights $\mathbf{w}^{(1)}, \dots, \mathbf{w}^{(k)}$ with zero gradients ($\mathbf{g} = \mathbf{0}$), freezing training permanently.

\subsection*{The Implementation Solution}

To resolve the type mutation and prevent the tracking graph from rupturing, a defensive type-enforcement guard is introduced mathematically and programmatically:
\[
\mathbf{x} = \begin{cases} \mathbf{x}, & \text{if } \mathbf{x} \text{ is an iterable collection} \\ \begin{bmatrix} \mathbf{x} \end{bmatrix}, & \text{if } \mathbf{x} \text{ is a naked scalar object} \end{cases}
\]
In code, this mathematical conditional check is executed directly within the layer forward pass:

\begin{lstlisting}[language=Python, style=pythonstyle]
class Layer:
    def __init__(self, nin, nout, activation='tanh'):
        self.neurons = [Neuron(nin, activation) for _ in range(nout)]

    def __call__(self, x):
        # TYPE GUARD: Force scalar objects back into an iterable tensor/list dimension
        if not isinstance(x, (list, tuple)):
            x = [x]
            
        outs = [n(x) for n in self.neurons]
        return outs[0] if len(outs) == 1 else outs
\end{lstlisting}

By guaranteeing that $x$ always preserves an explicit dimension sequence ($x \in \mathbb{R}^n$ where $n \ge 1$), the zip pairing mechanism is structurally preserved:
\[
\text{zip}\left(\mathbf{w}_j^{(k+1)}, \begin{bmatrix} a_1^{(k)} \end{bmatrix}\right) \implies \left(w_{j1}^{(k+1)}, a_1^{(k)}\right)
\]
The complete mathematical operation $w_{j1}^{(k+1)} \cdot a_1^{(k)}$ is registered by the autograd engine, preserving a valid localized Jacobian matrix and allowing gradient flow to propagate uninhibited through the entire network structure.

\end{document}